% Part: normal-modal-logic
% Chapter: axioms-systems
% Section: consistency

\documentclass[../../../include/open-logic-section]{subfiles}

\begin{document}

\olfileid{nml}{prf}{con}

\olsection{الاتساق}

% تصحيح مصدر (0440-I1): قُيّد الحكم الدلالي بالنماذج التي يصح فيها النسق.
ومن أهم ما ينظر فيه من أحوال مجموعات ال!!{formula}s اتساقها.
فمجموعة ال!!{formula}s غير متسقة متى كانت صيغة تناقض،
مثل~$\lfalse$، !!{derivable} منها؛ فإن امتنع ذلك كانت متسقة.
ولا تجتمع !!{formula}s المجموعة غير المتسقة على الصدق عند عالم واحد
في نموذج يصح فيه النسق. وأما العكس فهو مقصد مبرهنة التمام:
لكل مجموعة متسقة عالم تصدق عنده في نموذج نبنيه لهذا الغرض،
وهو «النموذج القانوني».

\begin{defn}
  تكون المجموعة~$\Gamma$ \emph{متسقة} بالنسبة إلى النسق~$\Sigma$، أو
  كما سنقول، \emph{$\Sigma$-متسقة}، إذا وفقط إذا كان
  $\Gamma \Proves/[\Sigma] \lfalse$.
\end{defn}

فالاعتبار في الاتساق بالنسق الذي يقع الاشتقاق فيه. ومن ذلك أن المجموعة
$\{\Box(p \lif q), \Box p, \lnot\Box q\}$ متسقة بالنسبة إلى منطق القضايا،
وليست~\Log{K}-متسقة. ومن ذلك أيضًا أن المجموعة
$\{\Diamond p, \Box\Diamond p \lif q, \lnot q\}$
ليست~\Log{K5}-متسقة.

\begin{prop}\ollabel{prop:consistencyfacts}
  لتكن~$\Gamma$ مجموعة من ال!!{formula}s. عندئذ:
  \begin{enumerate}
  \item تكون~$\Gamma$ $\Sigma$-متسقة إذا وفقط إذا وجدت
    !!{formula}~$!A$ بحيث~$\Gamma \Proves/[\Sigma] !A$.
  \item \ollabel{prop:consistencyfacts-b}%
    $\Gamma \Proves[\Sigma] !A$ إذا وفقط إذا لم تكن
    $\Gamma \cup \{\lnot!A\}$ $\Sigma$-متسقة.
  \item \ollabel{prop:consistencyfacts-c}%
    إذا كانت~$\Gamma$ $\Sigma$-متسقة، فلكل !!{formula}~$!A$ تكون إما
    $\Gamma \cup \{!A\}$ $\Sigma$-متسقة أو
    $\Gamma \cup \{\lnot!A\}$ $\Sigma$-متسقة.
  \end{enumerate}
\end{prop}

\begin{proof}
  يكفي في إثبات هذه الخصائص استعمال منطق القضايا الكلاسيكي.
  ولنبيّن حجة \olref{prop:consistencyfacts-c} بطريق المقابل.
  افرض أن كلا المجموعتين $\Gamma \cup \{!A\}$ و$\Gamma \cup \{\lnot!A\}$
  ليست $ \Sigma$-متسقة. فبحسب~\olref{prop:consistencyfacts-b} نحصل على
  $\Gamma, !A \Proves[\Sigma] \lfalse$ وعلى
  $\Gamma, \lnot !A \Proves[\Sigma] \lfalse$. فإذا رفعنا الفرض الإضافي
  بمبرهنة الاستنتاج كان
  $\Gamma \Proves[\Sigma] !A \lif \lfalse$ وكان
  $\Gamma \Proves[\Sigma] \lnot!A \lif \lfalse$. ولما كانت
  $(!A \lif \lfalse) \lif ((\lnot!A \lif \lfalse) \lif \lfalse)$
  مثال إحلال لقضية صادقة صدقًا تحليليًا، جاز تطبيق
  \olref[prp]{prop:derivabilityfacts}\olref[prp]{prop:derivabilityfacts-ruleT}
  فنستنتج $\Gamma \Proves[\Sigma] \lfalse$، وهو عدم اتساق المجموعة الأصلية.
\end{proof}

\end{document}
